
%(BEGIN_QUESTION)
% Copyright 2006, Tony R. Kuphaldt, released under the Creative Commons Attribution License (v 1.0)
% This means you may do almost anything with this work of mine, so long as you give me proper credit

Definer følgende kjemiske begreper: {\it løsning}, {\it løst stoff} (solute) og {\it løsemiddel} (solvent).

\underbar{file i00557}
%(END_QUESTION)





%(BEGIN_ANSWER)

En kjemisk {\it løsning} er en stabil og homogen blanding (på molekylært nivå) av to eller flere stoffer. {\it Solute} viser til stoffet/stoffene i mindre mengde, og {\it solvent} viser til hovedstoffet som solute er oppløst i.

\vskip 10pt

I mange tilfeller er det vilkårlig hva som kalles solute og solvent, som i en 50/50 alkohol-vann-løsning. Ofte er solute en gass eller et fast stoff, og solvent er en væske hvor gass(er) og/eller faststoff oppløses. Andre kombinasjoner er også mulige.

\vskip 10pt

Et eksempel der både solute og solvent er gass er luft. Her klassifiseres nitrogen som solvent fordi det er hovedkomponenten. Oksygen, karbondioksid, helium og de øvrige gassene blir da solute.

\vskip 10pt

Eksempler der solute er gass og solvent er væske: ammoniakk i vann, acetylen i aceton og nitrogen i blod.

\vskip 10pt

Eksempel der solute er gass og solvent er fast stoff: hydrogen oppløst i palladium.

\vskip 10pt

Eksempler der både solute og solvent er væsker: alkohol i vann og mange metallegeringer i smeltet tilstand.

\vskip 10pt

En klasse løsninger der solute er væske og solvent er fast stoff kalles {\it amalgamer}. Ett eksempel er kvikksølv (solute) oppløst i sølv (solvent), brukt i tannfyllinger.

\vskip 10pt

Eksempler der solute er fast stoff og solvent er væske: sukker i vann og kaliumklorid (KCl) i vann.

\vskip 10pt

Eksempler der både solute og solvent er faste stoffer: messing og bronse i fast tilstand. I disse legeringene er kobber solvent (hovedbestanddel), mens henholdsvis sink eller tinn er solute.

%(END_ANSWER)





%(BEGIN_NOTES)

%INDEX% Kjemi, grunnprinsipper: løsning, løst stoff og løsemiddel

%(END_NOTES)

